\chapter{Regularization and Variable Selection}
\label{cap:regularizacao}
% ── index ──────────────────────────────────────────────────────
\index{regularization|textbf}
\index{LASSO|textbf}
\index{lasso@\texttt{lasso()}|textbf}
\index{Ridge|textbf}
\index{ridge_reg@\texttt{ridge\_reg()}|textbf}
\index{Elastic Net|textbf}
\index{enet@\texttt{enet()}|textbf}
\index{sparsity|textbf}
\index{L1 penalty (lasso)}
\index{L2 penalty (ridge)}
\index{variable selection}

% ─────────────────────────────────────────────────────────────────────────────
\section{The Dimensionality Problem}

With $k$ regressors and small $N$, OLS suffers from \emph{overfitting}: it fits
the sample noise beyond the signal, increasing out-of-sample error. When
$k > N$, OLS is not even identified. Regularization penalizes the magnitude of
the coefficients, trading bias for reduced variance.

\subsection*{LASSO}

LASSO (Tibshirani 1996) minimizes:

\[
  \hat\beta^{\mathrm{LASSO}} = \arg\min_\beta \left\{ \sum(y_i - x_i'\beta)^2
  + \lambda \sum_{j=1}^k |\beta_j| \right\}
\]

The $\ell_1$ penalty forces small coefficients to exactly zero, performing
automatic \emph{variable selection}. The regularization path — LASSO for
different $\lambda$ — reveals which variables enter as $\lambda$ decreases
from $\infty$ (null model) to 0 (OLS).

\subsection*{Ridge}

The $\ell_2$ penalty of Ridge (Hoerl and Kennard 1970) shrinks coefficients
toward zero without zeroing them:

\[
  \hat\beta^{\mathrm{Ridge}} = (X'X + \lambda I)^{-1} X'y
\]

Useful when many variables have small effects (continuum of weak predictors).
Does not perform variable selection.

\subsection*{Elastic Net}

Elastic Net (Zou and Hastie 2005) combines $\ell_1$ and $\ell_2$:

\[
  \min_\beta \left\{ \sum(y_i - x_i'\beta)^2
  + \lambda \left[\rho\sum|\beta_j| + (1-\rho)\sum\beta_j^2\right] \right\}
\]

Inherits variable selection from LASSO and stability in correlated groups
from Ridge.

% ─────────────────────────────────────────────────────────────────────────────
\section{DGP Verification}

DGP with 10 regressors, only $x_1$ and $x_2$ with non-zero effect.
$N = 200$ — regime where regularization matters.

\begin{lstlisting}[caption={Sparse DGP — 2 of 10 active variables}]
seed(42)
// 10 regressors; y = 1 + 2*x1 + 3*x2 + 0*(x3..x10) + e
generate df x1  = rnormal()
...
generate df y   = 1.0 + 2.0*x1 + 3.0*x2 + e

let m_ols   = ols(y ~ x1+x2+x3+x4+x5+x6+x7+x8+x9+x10, df)
let m_lasso = lasso(y ~ x1+x2+x3+x4+x5+x6+x7+x8+x9+x10, df, alpha=0.1)
let m_ridge = ridge_reg(y ~ x1+x2+x3+x4+x5+x6+x7+x8+x9+x10, df, alpha=0.5)
let m_enet  = enet(y ~ x1+x2+x3+x4+x5+x6+x7+x8+x9+x10, df, alpha=0.5)
\end{lstlisting}

\subsection*{OLS}

\begin{lstlisting}[language={}, basicstyle=\ttfamily\small,
  frame=none, numbers=none, backgroundcolor=\color{codbg},
  xleftmargin=0pt, xrightmargin=0pt, breaklines=false]
OLS: R²=0.801  (good in-sample fit, but x3..x10 pollute with non-zero coefs)
const=0.68  x1=1.93***  x2=2.88***  x3=0.17  x4=0.07  ...  x10=-0.11
\end{lstlisting}

\subsection*{LASSO}

\begin{lstlisting}[language={}, basicstyle=\ttfamily\small,
  frame=none, numbers=none, backgroundcolor=\color{codbg},
  xleftmargin=0pt, xrightmargin=0pt, breaklines=false]
LASSO α=0.1: R²=0.772  active vars: 2
  x1=1.552  x2=2.506  x3..x10 = 0.000  (correctly zeroed)
\end{lstlisting}

With $\alpha = 0{,}1$ LASSO identifies the two true predictors and
zeros the eight irrelevant ones.

\subsection*{Ridge}

\begin{lstlisting}[language={}, basicstyle=\ttfamily\small,
  frame=none, numbers=none, backgroundcolor=\color{codbg},
  xleftmargin=0pt, xrightmargin=0pt, breaklines=false]
Ridge α=0.5: R²=0.801
  x1=1.878  x2=2.793  x3=0.172  x4=0.081  ...
  (shrinks but zeros none)
\end{lstlisting}

\subsection*{Elastic Net}

\begin{lstlisting}[language={}, basicstyle=\ttfamily\small,
  frame=none, numbers=none, backgroundcolor=\color{codbg},
  xleftmargin=0pt, xrightmargin=0pt, breaklines=false]
ElasticNet α=0.5 l1_ratio=0.5: R²=0.589  active vars: 2
  x1=0.789  x2=1.550  x3..x10 = 0.000
\end{lstlisting}

Elastic Net also selects $x_1$ and $x_2$, but with greater total penalization
(combining $\ell_1$ and $\ell_2$) the coefficients are more shrunk.

% ─────────────────────────────────────────────────────────────────────────────
\section{Syntax}

\begin{lstlisting}
lasso(y ~ x1 + x2 + x3, df, alpha=0.1)      // alpha = penalization lambda
ridge_reg(y ~ x1 + x2 + x3, df, alpha=0.5)
enet(y ~ x1 + x2 + x3, df, alpha=0.5)       // alpha = L1/(L1+L2) mix
\end{lstlisting}

\begin{center}
\begin{tabular}{lll}
\toprule
\textbf{Parameter} & \textbf{LASSO} & \textbf{Ridge / Enet} \\
\midrule
\texttt{alpha} & penalty $\lambda$ & penalty $\lambda$ \\
\texttt{alpha=0} & OLS & pure Ridge (Enet) \\
\texttt{alpha=1} & max L1 & pure L1 (Enet) \\
\bottomrule
\end{tabular}
\end{center}

\begin{nota}
  The optimal value of $\alpha$ is typically chosen by cross-validation
  ($k$-fold). In small samples, LASSO with $\alpha$ via \emph{BIC-LASSO}
  tends to be more conservative than via CV.
\end{nota}
